\section{Des collections d'exception au service du droit}

La bibliothèque Cujas regorge d'une richesse documentaire géré par le Département du Patrimoine, de la conservation et de la numérisation (DPCN) qui gère \enquote{40 0000 documents patrimoniaux du XVe au XXe siècle conservés en salle du Patrimoine}. Elle est forte d'un des plus importants fonds juridiques à l'échelle internationale.

La richesse des collections de Cujas repose historiquement sur des fonds spécialisés et des ouvrages rares. L'établissement est ainsi internationalement reconnu pour conserver un \enquote{ensemble unique de ressources, de collections en sciences juridiques, économiques et politiques}. Grâce à cette spécialisation, la bibliothèque est devenue le pôle de référence incontournable pour les chercheurs et les professionnels du droit.

Au-delà de ses monographies, la bibbliothèque abrite un précieux gisement de sources primaires comprenant \enquote{les thèses, les mémoires et les publications officielles de la Société des Nations, des Nations Unies, du Journal officiel France et des anciennes colonies}. À cet inventaire s'ajoute une remarquable \enquote{collection de travaux académiques imprimés en Allemagne [du] XVIe [aux] XIXe siècle}, constituve du fonds de Jean-Marie Pardessus. Rédigés pour la plupart en latin, ces travaux retracent avec une grande précision la sociabilité universitaire et l'histoire de l'enseignement de l'ancienne Faculté de droit. 

Par ailleurs, la bibliothèque Cujas s'est particulièrement efforcée de constituer et de préserver \enquote{une collection quasi exhaustive de thèses de droit soutenues à la Faculté de droit}. Aujourd'hui encore, l'établissement poursuit activement cette mission patrimoniale de collecte pour les thèses de droit, d'économie et de science politique soutenues au sein des universités Panthéon-Sorbonne (Paris I), et Panthéon-Assas (Paris II). Le support de ces travaux varie selon la période, passant de la forme imprimée, à la forme microfichée ou numérique.

Les collections de la bibliothèque Cujas s'étendent bien au-delà des monographies traditionnelles pour intégrer des ressources uniques, telles que les dissertations d'agrégation de droit, les manuels de pédagogie du droit. Les manuels de pédagogie du droit permettent aux chercheurs d'observer et de constater les évolutions du droit français et de son enseignement tout au long du XIXe siècle. C'est précisément pour valoriser ce dernier corpus qu'a été conçu le projet HOPPE-Droit (Hypothèses d'Observation des Productions Pédagogiques Editées en Droit), financé par le groupe scientifique CollEx-Persée. Grâce à sa richesse patrimoniale, la bibliothèque Cujas a pu s'appuyer sur cette dynamique collective pour \enquote{sonder l'état de ses collections et enrichir ses données, tout en offrant aux chercheurs un instrument de travail inédit}. 

Sur le plan méthodologique, HOPPE-Droit se définit comme un outil numérique d'analyse relationnelle conçu pour étudier l'évolution de l'enseignement du droit français du XIXe au XXe siècle. La particularité de cet outil réside dans sa capacité à croiser \enquote{des éléments des notices bibliographiques de chacun des ouvrages retenus dans le corpus qu'elle fait résonner les uns par rapport aux autres (titres, auteurs, domaine, éditeurs, collections)}. A travers la mise en relation par des liens -- telles que les tittres, les auteurs, les domaines disciplinaires, les éditeurs et les collections --, l'outil permet de sortir de simples descriptions isolées pour cartographier le squelette conceptuel de la littérature pédagogique d'Ancien Régime et de l'époque contemporaine.

Cette modélisation de données massives offre de nouvelles perspectives de recherche, permettant de mettre en évidence les \enquote{tendances historiques qui se sont faites jour en matières d'ouvrages juridiques à vocation pédagogique} et de cartographier \enquote{les réseaux d'édition qui structurent le champ du droit et d'en suivre l'évolution dans le temps}. Finalement, cette entreprise scientifique n'a pu voir le jour que grâce à une collaboration étroite et les professionnels des bibliothèques. Ensemble ils partagent un objectif scientifique commun : \enquote{suivre les évolutions de l'enseignement du droit, saisies à travers le prisme de l'édition juridique à vocation pédagogique}. 
	
Les fonds spécialisés de la bibliothèque Cujas se caractérisent par leur remarquable diversité et se répartissent en six grands ensembles documentaires distincts. Ces collections comprennent les publications officielles de la Société des Nations (SDN), celles des Nations Unies (ONU), celle de la collection historique du Journal Officiel (JO), le fonds universitaires de Jean-Marie Pardessus, ainsi qu'un ensemble de cartes postales et les archives historiques de la Faculté de droit de Paris. 

Le premier pôle regroupe les sources officielles et étatiques. La collection des publications officielles de la SDN a été constitué dès les années 1920, couvrant l'intégralité de sa production éditoriale entre 1919 et 1946 à travers 534 volumes rédigés en français et en anglais. Pour l'ONU, la bibliothèque Cujas bénéficie du statut de dépositaire depuis 1949, accumulant un fonds colossal de près de 3 800 cartons, dont l'approvisionnement physique a cessé en 2010 à la suite de la dématérialisation des documents officiels. Enfin, la bibliothèque conserve l'intégralité du Journal Officiel sur une période s'étendant de 1870 à 2015. Ces numéros historiques sont consultables sous leur forme imprimée jusqu'en 2015, tandis que leur version numérique est accessible en ligne pour les éditions postérieures à 2004.

Le fonds Jean-Marie Pardessus (1772-1853) constitue quant à lui un monument de l'histoire universitaire européenne. Cette collection de près de 9 000 travaux académiques en droit couvre une vaste période allant du XVIe au XIXe siècle. Ces thèses et dissertations ont la particularité d'avoir \enquote{principalement été rédigés en latin et soutenues dans les universités allemandes}. Ce fonds ancien est entré dans les collections de la bibliothèque Cujas par l'intermédiaire d'un don consenti par M. de Rozière, petit-fils de Jean-Marie Pardessus, à la Faculté de droit de Paris. 

Enfin, la mémoire de la Faculté de droit de Paris fait l'objet d'une double valorisation, à la fois iconographique et archivistique. D'une part, la bibliothèque conserve un ensemble original de cartes postales \enquote{caricaturales représentant des professeurs de la Faculté de droit de Paris au début du XXe siècle}, couvrant la période de 1900 à 1914. D'autre part, le fonds des archives de la Faculté de droit de Paris retrace \enquote{la vie institutionnelle, pédagogique et politique de la Faculté de droit (XIXe-XXe siècle)}. Composé de sources administratives, de papiers personnels de professeurs et de manuscrits, ce derniers fonds ofre un témoignage précieux sur l'évolution de l'enseignement juridique parisien.
	
L'exhaustivité de ses fonds documentaires, qui fait de l'établissement une bibliothèque d'excellence et de référence mondiale dans le domaine juridique, ne peut se cantonner à un rôle de conservation passive. C'est précisément cette richesse et cette rareté qui en font une institution majeure au sein du paysage universitaire et qui justifient son intégration dans les réseaux de recherche.